% =====================================================================
%  On the Necessary Substructures of Finite Contact Graphs
%  A conditional-existence account of why a finite world is structured.
% =====================================================================
\documentclass[11pt,a4paper]{article}

\usepackage[utf8]{inputenc}
\usepackage[T1]{fontenc}
\usepackage{lmodern}
\usepackage[a4paper,margin=2.5cm]{geometry}
\usepackage{amsmath,amssymb,amsthm,mathtools}
\usepackage{bm}
\usepackage{mathrsfs}
\usepackage{booktabs}
\usepackage{array}
\usepackage{enumitem}
\usepackage{microtype}
\usepackage{graphicx}
\usepackage[round,sort&compress]{natbib}
\usepackage{tikz}
\usetikzlibrary{arrows.meta,positioning,calc,decorations.pathmorphing,fit}
\usepackage[colorlinks=true,linkcolor=blue!45!black,
            citecolor=blue!45!black,urlcolor=blue!45!black]{hyperref}
\usepackage{cleveref}

% ---- Theorem environments ----
\newtheoremstyle{thmstyle}{\topsep}{\topsep}{\itshape}{}{\bfseries}{.}{.5em}{}
\theoremstyle{thmstyle}
\newtheorem{theorem}{Theorem}[section]
\newtheorem{lemma}[theorem]{Lemma}
\newtheorem{proposition}[theorem]{Proposition}
\newtheorem{corollary}[theorem]{Corollary}
\theoremstyle{definition}
\newtheorem{definition}[theorem]{Definition}
\newtheorem{axiom}{Axiom}
\newtheorem{principle}[theorem]{Principle}
\newtheorem{construction}[theorem]{Construction}
\newtheorem{example}[theorem]{Example}
\theoremstyle{remark}
\newtheorem{remark}[theorem]{Remark}
\newtheorem{notation}[theorem]{Notation}

\crefname{axiom}{Premise}{Premises}
\Crefname{axiom}{Premise}{Premises}
\crefname{principle}{Principle}{Principles}
\Crefname{principle}{Principle}{Principles}

% ---- Shorthand ----
\newcommand{\R}{\mathbb{R}}
\newcommand{\N}{\mathbb{N}}
\newcommand{\Z}{\mathbb{Z}}
\newcommand{\Gph}{\Gamma}                 % the contact graph
\newcommand{\V}{V}                        % vertex set
\newcommand{\E}{E}                        % edge set
\newcommand{\wt}{w}                       % edge weight
\newcommand{\floorc}{\beta}               % the floor
\newcommand{\Nbr}{N}                      % neighbourhood
\newcommand{\co}{\complement}             % complement
\newcommand{\cut}{\partial}               % edge boundary / cut
\newcommand{\Walk}{W}                     % a walk
\newcommand{\gate}{g}                     % gating function
\newcommand{\Agent}{A}                    % agent subgraph
\newcommand{\Res}{\rho}                   % residual invariant
\newcommand{\Quot}{Q}                     % quotient graph
\newcommand{\rec}{M}                      % committed record
\newcommand{\Ver}{\mathsf{V}}             % verifier
\newcommand{\abs}[1]{\lvert#1\rvert}
\newcommand{\norm}[1]{\lVert#1\rVert}
\newcommand{\restr}{\!\restriction\!}

\title{\textbf{On the Necessary Substructures of Finite Contact Graphs:\\[2pt]
A Conditional-Existence Account of Why a Finite World Is Structured}}

\author{%
Kundai Farai Sachikonye\\
\small Technical University of Munich\\
\small AIMe Registry for Artificial Intelligence\\
\small \texttt{kundai.sachikonye@wzw.tum.de}}

\date{\today}

\begin{document}
\maketitle

% =====================================================================
\begin{abstract}
\noindent
We ask a structural question---\emph{why is a world structured at all?}---and
answer it with a single hypothesis about graphs. A \emph{contact graph} is a
finite weighted graph whose vertices are the parts that have been individuated
within a whole and whose edges carry the cost of separating one part from
another. We take as our only premise that such a graph is \emph{finite} and
that its edge weights are bounded below by a positive constant; we call this the
Finiteness Premise. From it alone we derive, as a chain of existence theorems, a
ladder of substructures that any finite contact graph must contain. We prove:
\textbf{(T0)} a strictly positive \emph{floor} on every separating boundary---no
two parts are adjacent at zero cost; \textbf{(T1)} that each vertex is fixed only
by its complement-neighbourhood, so that individuation must proceed by negation
and admits no selector without regress; \textbf{(T2)} that the
\emph{cut-residual}---the inter-part boundary content---is bounded below and is
invariant under every relabelling (isomorphism) of the graph, while no single
vertex realises it, so the residual is a region and never a point; \textbf{(T3)}
that this residual is the conserved quantity under all re-encodings of the
graph's vertices, the unique invariant from which the relabellable surface
derives; \textbf{(T4)} that every connected subgraph carries its own such
invariant under its automorphisms; \textbf{(T5)} that selecting a walk on the
graph requires an edge-gating function, without which the next step is
undetermined; \textbf{(T6)} that the committed-edge count along any walk is
monotone, so revisiting a vertex is never returning to a state and resemblance
never entails identity; \textbf{(T7)} that a coherent family of subgraphs acting
as one admits a single gating function on their quotient graph, and only one;
and \textbf{(T8)} that, because the subgraph invariant is complete, contact
graphs are specifiable and hence instantiable---authored finite structures
exist---subject to three bounds we prove: their walks are not foreseeable from
their specification, an authored copy is a distinct structure, and authorship is
limited to the resolution of the floor. We close with a Master Theorem: every
substructure T0--T8 is forced for finite graphs and is \emph{not} forced for
infinite or zero-floor graphs. Structure is therefore neither imposed on a world
nor intrinsic to it independent of its parts; it is the necessary form a world
must take to be resolved by finitely many bounded parts. \emph{Structure is the
shadow of finitude.} Every claim is proved from the single premise; uncertain
material has been omitted, and the interpretations a reader may attach are
confined to clearly marked remarks on which no theorem depends.
\end{abstract}

\noindent\textbf{Keywords:} contact graph; finite weighted graph; edge-weight
floor; individuation by negation; cut invariant; graph isomorphism invariant;
gated walk; monotone walk; quotient graph; conditional existence; necessary
substructure; authored structure.

\tableofcontents

% =====================================================================
\section{Introduction}
\label{sec:intro}
% =====================================================================

\subsection{The question}

Why is a world structured at all? One might expect a world to be either a
featureless continuum, in which nothing is marked off from anything else, or an
arbitrary heap, in which everything is marked off from everything else at no
cost. Neither is what is observed: parts are distinguished, but their
distinction is bounded---there is a cost to telling one part from another, and
the parts compose into a definite architecture. We argue that this architecture
is not contingent. It is forced, in full, by a single fact: that the parts are
\emph{finite in number} and \emph{finitely separated}. The structure of a world
is the shadow its finiteness casts.

We make this precise in the language of graphs, because every relation of
``this part, separated from that part'' is natively an edge between two
vertices, and any architecture of parts is some graph. The object of the paper
is therefore a single graph---a \emph{contact graph}---whose vertices are the
individuated parts of a whole and whose weighted edges record the cost of the
separations between them. We pose one hypothesis about this graph and derive
everything else.

\subsection{Method: conditional existence}

The paper proves no empirical claim. It establishes a conditional: \emph{if} a
finite contact graph exists, \emph{then} a determinate ladder of substructures
must exist within it. Each rung is an existence theorem of the form ``under the
Finiteness Premise, the following is forced.'' We assume nothing about what the
vertices \emph{are}; we assume only that they are finite in number and
positively separated, and we read off what their graph cannot lack. The
deductions are combinatorial and order-theoretic throughout. Where a tempting
claim could not be proved at this standard it has been removed rather than
weakened.

\subsection{The single premise}

\begin{axiom}[Finiteness Premise]
\label{ax:finite}
A \emph{contact graph} is a finite weighted graph $\Gph=(\V,\E,\wt)$ with
$\abs{\V}<\infty$, in which every edge weight is bounded below by a positive
constant: there is $\floorc>0$ with $\wt(e)\ge\floorc$ for all $e\in\E$. We
assume throughout that a contact graph exists and that it is connected.
\end{axiom}

The premise is almost nothing: a finite graph whose edges cost something. Its
consequences are almost everything we mean by structure.

\subsection{The ladder of necessary substructures}

Under \Cref{ax:finite} we prove the following chain, each rung from the premise
together with those above it.

\begin{enumerate}[label=\textbf{(T\arabic*)},leftmargin=3em,start=0]
\item \emph{The floor exists} (\Cref{sec:floor}): every separating boundary has
weight $\ge\floorc>0$; no two parts are adjacent at zero cost.
\item \emph{Individuation is by negation} (\Cref{sec:negation}): a vertex is
fixed only as the complement of its non-neighbours; no selector fixes it without
regress.
\item \emph{Truth is a residual region} (\Cref{sec:residual}): the cut-residual
is bounded below and realised by no single vertex, hence a region, never a
point.
\item \emph{The residual is the conserved invariant} (\Cref{sec:residual}):
under every relabelling of vertices the residual is preserved; it is the unique
invariant from which the relabellable surface derives.
\item \emph{Each subgraph has a private invariant} (\Cref{sec:invariant}): every
connected subgraph carries its own invariant under its automorphisms.
\item \emph{Selection requires a gate} (\Cref{sec:gating}): a walk is
undetermined without an edge-gating function that biases the next step.
\item \emph{Histories do not return} (\Cref{sec:monotone}): the committed-edge
count is monotone; revisiting a vertex is not returning to a state.
\item \emph{Plurality forces a single quotient gate} (\Cref{sec:quotient}): a
coherent family of subgraphs acting as one admits exactly one gating function on
their quotient.
\item \emph{Authored structures exist, and are bounded} (\Cref{sec:construct}):
the subgraph invariant is complete, so contact graphs are specifiable and
instantiable; their walks are unforeseeable from the specification, an authored
copy is a distinct structure, and authorship reaches only to the floor.
\end{enumerate}

\Cref{sec:master} proves the Master Theorem: T0--T8 are forced for finite
graphs and not forced when finiteness or the positive floor is dropped.

\subsection{On the word ``contact''}

We use \emph{contact} in one operational sense: two parts are \emph{in contact}
when the graph carries an edge between them, and the weight of that edge is the
cost of the separation that individuates each from the other. Nothing in the
paper depends on any informal reading of the word; ``contact'' names an edge and
its weight, and the theorems about contact graphs are theorems about finite
weighted graphs in exactly this sense.

\subsection{Why finiteness is the source}

Finiteness is not a convenience of presentation but the origin of every
structure below. On an infinite graph a vertex may be separated from the rest by
edges of vanishing weight, the residual may be driven to zero, and a walk may be
specified to unbounded depth at fixed cost; none of the substructures is then
forced. It is the finitude of the vertex set, together with the positive floor
on edge weights, that forbids the zero-cost separation and compels the ladder.
We make the premise exactly strong enough to force the structures and no
stronger, and we mark, at each rung, where finiteness is used
(\Cref{sec:master}).

\subsection{Relation to existing work}

The paper is pure mathematics. It uses graph theory~\citep{bondymurty2008,%
diestel2017,west2001,bollobas1998}, algebraic and spectral graph
theory~\citep{godsilroyle2001,chung1997,fiedler1973,cheeger1970}, the theory of
walks and Markov chains on graphs~\citep{lovasz1993,levinperes2017}, order and
fixed-point theory~\citep{daveypriestley2002,birkhoff1967,tarski1955},
enumerative combinatorics and partitions~\citep{stanley2012,andrews1976},
self-reference and diagonal arguments~\citep{russell1903,cantor1891,lawvere1969,%
godel1931}, and information theory~\citep{shannon1948,cover2006}. Where the
development meets the sorites and the logic of vagueness we note the
connection~\citep{black1937,williamson1994} without taking any position in that
literature; our treatment is geometric, not a logic of partial truth. A
reasonable reader will recognise that the abstract premise is satisfied by a
range of concrete systems of finitely many bounded parts; we make no such
identification and prove everything of the abstract graph alone. Such
interpretive readings as we offer are confined to remarks labelled accordingly,
on which no theorem depends.

\subsection{Organisation}

\Cref{sec:setting} fixes the contact graph and its derived objects.
\Cref{sec:floor} proves T0. \Cref{sec:negation} proves T1. \Cref{sec:residual}
proves T2 and T3. \Cref{sec:invariant} proves T4. \Cref{sec:gating} proves T5.
\Cref{sec:monotone} proves T6. \Cref{sec:quotient} proves T7.
\Cref{sec:construct} proves T8. \Cref{sec:master} proves the Master Theorem and
concludes.


% =====================================================================
\section{The contact graph and its derived objects}
\label{sec:setting}
% =====================================================================

We fix notation and the objects on which the ladder is built. Throughout,
$\Gph=(\V,\E,\wt)$ is a contact graph in the sense of \Cref{ax:finite}: a finite,
connected, weighted graph with weight floor $\floorc>0$. Graphs are simple
(no loops or multi-edges) unless stated; all standard graph-theoretic terms are
as in~\citet{bondymurty2008,diestel2017}.

\subsection{Vertices, edges, weights}

\begin{definition}[Contact graph]
\label{def:contactgraph}
A \emph{contact graph} is a triple $\Gph=(\V,\E,\wt)$ where $\V$ is a finite set
of \emph{vertices} (the individuated parts), $\E\subseteq\binom{\V}{2}$ is a set
of \emph{edges}, and $\wt\colon\E\to[\floorc,\infty)$ is a \emph{weight}
assigning to each edge the cost of the separation it records, with
$\floorc>0$ the \emph{floor}. For $S,T\subseteq\V$ disjoint, the
\emph{cut} between them is
\[
\cut(S,T)\;=\;\{\,uv\in\E : u\in S,\ v\in T\,\},
\qquad
\wt(\cut(S,T))\;=\!\!\sum_{e\in\cut(S,T)}\!\!\wt(e).
\]
We write $\cut(S):=\cut(S,\V\setminus S)$ for the boundary of $S$.
\end{definition}

\begin{definition}[Neighbourhood and complement-neighbourhood]
\label{def:nbr}
For $v\in\V$, the \emph{neighbourhood} is
$\Nbr(v)=\{u\in\V : uv\in\E\}$ and the \emph{closed neighbourhood} is
$\Nbr[v]=\Nbr(v)\cup\{v\}$. The \emph{complement-neighbourhood} is
\[
\co\Nbr[v]\;:=\;\V\setminus\Nbr[v],
\]
the set of vertices that $v$ neither is nor touches---``everything $v$ is not.''
\end{definition}

\begin{remark}[The floor is the whole content of ``contact'']
\label{rem:floor-content}
A contact graph differs from an arbitrary weighted graph in exactly one
respect: the positive lower bound $\floorc$ on edge weights. This single
constraint---no separation is free---is what every theorem below turns on. A
weighted graph permitting $\wt(e)\to0$ is not a contact graph; it is the
degenerate object against which our results are stated by contrast
(\Cref{sec:master}).
\end{remark}

\subsection{Subgraphs as parts-within-the-whole}

\begin{definition}[Induced subgraph; agent]
\label{def:subgraph}
For $\varnothing\neq U\subseteq\V$, the \emph{induced subgraph} $\Gph[U]$ has
vertex set $U$ and edge set $\{uv\in\E : u,v\in U\}$ with inherited weights. A
connected induced subgraph $\Agent=\Gph[U]$ with $1<\abs{U}<\abs{\V}$ is called a
\emph{part} (or, where we wish to emphasise that a part may itself individuate
others, an \emph{agent}); its \emph{boundary} is the cut $\cut(U)$.
\end{definition}

\begin{remark}[Vocabulary]
\label{rem:agent-word}
The word ``agent'' is used only as a name for a connected induced subgraph that
participates in further separations; it carries no further commitment. Every
theorem stated for an agent is a theorem about an induced subgraph and its
boundary cut.
\end{remark}

\subsection{Isomorphism, automorphism, invariant}

\begin{definition}[Weighted isomorphism]
\label{def:iso}
A \emph{weighted isomorphism} $\phi\colon\Gph\to\Gph'$ is a bijection
$\phi\colon\V\to\V'$ with $uv\in\E\iff\phi(u)\phi(v)\in\E'$ and
$\wt'(\phi(u)\phi(v))=\wt(uv)$. An \emph{automorphism} is a weighted isomorphism
$\Gph\to\Gph$; the automorphisms form the group $\mathrm{Aut}(\Gph)$. A
\emph{graph invariant} is a function $\iota$ on contact graphs with
$\iota(\Gph)=\iota(\Gph')$ whenever $\Gph\cong\Gph'$.
\end{definition}

\begin{remark}[Re-encoding is relabelling]
\label{rem:reencode}
A change of representation that preserves which parts are separated from which,
and at what cost, is precisely a weighted isomorphism: it relabels vertices
without altering the graph. Quantities that survive all such relabellings are
the graph invariants; quantities tied to particular labels do not. This
distinction---surface labels versus invariants---is the form in which
``what is conserved under re-encoding'' appears below (\Cref{sec:residual}).
\end{remark}

\subsection{Walks and gating}

\begin{definition}[Walk]
\label{def:walk}
A \emph{walk} of length $\ell$ is a sequence
$\Walk=(v_0,v_1,\dots,v_\ell)$ with $v_{i}v_{i+1}\in\E$ for each $i$. A walk may
repeat vertices and edges. The \emph{committed multiset} of $\Walk$ is the
multiset of edges $\{v_0v_1,\dots,v_{\ell-1}v_\ell\}$ traversed, and its
\emph{committed count} is $\rec(\Walk)=\ell$, the number of steps taken.
\end{definition}

\begin{definition}[Gating function]
\label{def:gate}
A \emph{gating function} on $\Gph$ is a map
$\gate\colon\V\times\E\to[0,1]$ assigning to each vertex $v$ and incident edge
$e\ni v$ a forward weight $\gate(v,e)$, with $\sum_{e\ni v}\gate(v,e)\le 1$ for
each $v$. A walk is \emph{$\gate$-admissible} if every step $v_iv_{i+1}$ has
$\gate(v_i,v_iv_{i+1})>0$. The gate \emph{selects} a step when it assigns
positive forward weight to exactly one incident edge.
\end{definition}

\begin{remark}[Why gating is needed at all]
\label{rem:gate-need}
A contact graph is, in general, far from a path: a vertex has many incident
edges, so a walk's next step is underdetermined by the graph alone. Some
additional datum---a gating function---is required to pick the next edge. The
necessity of such a datum, not its particular form, is what
\Cref{sec:gating} establishes. We treat $\gate$ abstractly; any rule that biases
the next step is a gating function in the sense of \Cref{def:gate}.
\end{remark}

\subsection{The quotient of a partition}

\begin{definition}[Quotient graph]
\label{def:quotient}
Let $\mathcal{P}=\{U_1,\dots,U_k\}$ be a partition of $\V$ into connected
classes. The \emph{quotient graph} $\Gph/\mathcal{P}=(\Quot,\E_{\Quot},\wt_{\Quot})$
has vertex set $\Quot=\{U_1,\dots,U_k\}$, an edge $U_iU_j$ whenever
$\cut(U_i,U_j)\neq\varnothing$, and weight
$\wt_{\Quot}(U_iU_j)=\wt(\cut(U_i,U_j))$. The quotient is itself a weighted
graph; when each class is a part in the sense of \Cref{def:subgraph}, the
quotient records the contacts \emph{between} parts.
\end{definition}

\begin{remark}[A graph of parts is again a graph]
\label{rem:quotient}
\Cref{def:quotient} lets the same theory apply one level up: the parts of a
contact graph, taken as vertices and joined by their mutual cuts, form a
weighted graph to which \Cref{ax:finite} again applies (its floor is the
minimum nonzero cut weight). This self-similarity is used in
\Cref{sec:quotient}.
\end{remark}

\begin{lemma}[The quotient inherits a floor]
\label{lem:quotient-floor}
If $\Gph$ is a contact graph with floor $\floorc$ and $\mathcal{P}$ is a
partition into connected classes such that every pair of classes joined in
$\Gph/\mathcal{P}$ is joined by at least one edge of $\Gph$, then
$\Gph/\mathcal{P}$ is a contact graph with floor $\ge\floorc$.
\end{lemma}

\begin{proof}
Each edge $U_iU_j$ of the quotient has weight
$\wt_{\Quot}(U_iU_j)=\wt(\cut(U_i,U_j))=\sum_{e\in\cut(U_i,U_j)}\wt(e)$, a sum of
at least one term each $\ge\floorc$ by \Cref{ax:finite}. Hence
$\wt_{\Quot}(U_iU_j)\ge\floorc>0$. The quotient is finite (at most $\abs{\V}$
classes) and, since $\Gph$ is connected, connected. Thus it satisfies
\Cref{ax:finite}.
\end{proof}


% =====================================================================
\section{T0: the floor}
\label{sec:floor}
% =====================================================================

The first substructure is the floor itself: in a contact graph no part can be
separated from the rest at zero cost. The statement is immediate from the
premise, but its consequences---restated for cuts, for parts, and for the
limiting behaviour under refinement---are what the later rungs use, so we record
them with care.

\subsection{Every separation is costly}

\begin{theorem}[Floor Theorem]
\label{thm:floor}
Let $\Gph$ be a contact graph with floor $\floorc>0$. Then every nonempty proper
part $U\subsetneq\V$ has boundary weight bounded below by the floor:
\[
\boxed{\ \wt(\cut(U))\ \ge\ \floorc\ >\ 0.\ }
\]
In particular no part is separated from its complement by a cut of zero
weight: a \emph{sharp} separation, $\wt(\cut(U))=0$, is impossible.
\end{theorem}

\begin{proof}
Since $\Gph$ is connected (\Cref{ax:finite}) and $U$ is a nonempty proper
subset, there is at least one edge with one endpoint in $U$ and the other in
$\V\setminus U$; otherwise $U$ and $\V\setminus U$ would be unjoined, splitting
the connected graph $\Gph$ into two nonempty components, a contradiction. Hence
$\cut(U)\neq\varnothing$, and
\[
\wt(\cut(U))\;=\!\!\sum_{e\in\cut(U)}\!\!\wt(e)\;\ge\;\wt(e^\ast)\;\ge\;\floorc>0
\]
for any $e^\ast\in\cut(U)$, using $\wt(e)\ge\floorc$ from \Cref{ax:finite}. A
sharp separation would require $\wt(\cut(U))=0<\floorc$, impossible.
\end{proof}

\begin{definition}[Boundary cost; the realised floor]
\label{def:boundarycost}
The \emph{boundary cost} of a part $U$ is $b(U):=\wt(\cut(U))$. The
\emph{realised floor} of $\Gph$ is
\[
\floorc_\Gph\;:=\;\min_{\varnothing\neq U\subsetneq\V} b(U)\;\ge\;\floorc>0,
\]
the least boundary cost of any part. By \Cref{thm:floor} it is strictly
positive, and by finiteness of $\V$ the minimum is attained.
\end{definition}

\begin{corollary}[No zero-cost individuation]
\label{cor:no-zero}
In a contact graph, individuating any part from the rest deposits boundary cost
at least $\floorc_\Gph>0$. Equivalently, the map $U\mapsto b(U)$ is bounded
away from zero on nonempty proper parts.
\end{corollary}

\begin{proof}
Immediate from \Cref{thm:floor} and \Cref{def:boundarycost}.
\end{proof}

\subsection{Refinement lowers but never abolishes the floor}

A finer description splits parts into smaller ones. We show this can only lower
the realised floor toward, but never to, zero---so long as the graph remains
finite with a positive weight floor.

\begin{proposition}[Refinement bound]
\label{prop:refine}
Let $\Gph'$ be obtained from $\Gph$ by any sequence of vertex splits that
preserves the floor $\floorc$ on all edges \textup{(}each new edge again has
weight $\ge\floorc$\textup{)}. Then the realised floor satisfies
$\floorc\le\floorc_{\Gph'}$, and in particular $\floorc_{\Gph'}>0$. The realised
floor can decrease under refinement but is bounded below by $\floorc$ and never
reaches zero at any finite stage.
\end{proposition}

\begin{proof}
$\Gph'$ is again finite (a finite number of splits of a finite graph) and, by
hypothesis, has all edge weights $\ge\floorc$, hence is a contact graph with
floor $\floorc$. By \Cref{thm:floor} applied to $\Gph'$, every part of $\Gph'$
has boundary cost $\ge\floorc$, so $\floorc_{\Gph'}\ge\floorc>0$. Splitting can
only introduce parts with smaller boundary cost, so $\floorc_{\Gph'}\le
\floorc_\Gph$ is possible; in all cases the bound $\floorc_{\Gph'}\ge\floorc$
holds. Only in the non-admissible limit of infinitely many splits, or of edge
weights tending to $0$, would the bound vanish---both excluded by
\Cref{ax:finite}.
\end{proof}

\begin{remark}[What T0 does and does not assert]
\label{rem:floor-scope}
\Cref{thm:floor} does not assert that some ambient continuum has thick
boundaries; it asserts that a \emph{finite, positively weighted} graph cannot
exhibit a zero-cost cut. The floor is a property of the graph-with-its-weights,
i.e. of the finite system of separated parts, not of any underlying space. This
is the precise sense in which, for finitely many bounded parts, ``every
distinction costs something.'' Every later rung consumes \Cref{thm:floor} as the
guarantee that boundaries, residuals, and gated steps are bounded away from
zero.
\end{remark}


% =====================================================================
\section{T1: individuation by negation}
\label{sec:negation}
% =====================================================================

The floor says separations cost something; it does not yet say how a part is
\emph{determined}. We now show that in a contact graph with no distinguished
vertex, a part can be fixed only by its complement---by everything it is
not---and that any attempt to fix it positively, by naming a distinguished
element, regresses. This is the graph-theoretic form of individuation by
negation.

\subsection{No distinguished vertex}

\begin{definition}[Selector; selector-free determination]
\label{def:selector}
A \emph{selector} for a part $U\subseteq\V$ is a rule that fixes $U$ by naming a
distinguished vertex of $U$ in advance of the separation. A determination of $U$
is \emph{selector-free} if it names no distinguished vertex of $U$. We say
$\Gph$ has \emph{no distinguished vertex} if its weighted automorphism group
acts with no fixed datum singling out a vertex in advance---formally, if the
specification of $\Gph$ provided to a determining rule is invariant under
$\mathrm{Aut}(\Gph)$, so that no vertex is named by the data alone.
\end{definition}

\begin{remark}[The setting of T1]
\label{rem:nodist}
T1 concerns the determination of parts from the graph's own data, with no
external label supplied. This is the relevant case for a self-contained finite
system: there is no oracle outside $\Gph$ to point at a vertex, so any vertex
named must be named \emph{by the structure}, and the structure is fixed only up
to $\mathrm{Aut}(\Gph)$. We make this precise and show negation is the only
route.
\end{remark}

\subsection{Complementation determines a part}

\begin{definition}[Negation-set]
\label{def:negationset}
For a part $U\subseteq\V$, its \emph{negation-set} is the complement
$\co U:=\V\setminus U$. A part $U$ is \emph{individuated by negation} if it is
determined as the unique part whose negation-set is a given $\co U$; that is,
$U=\V\setminus\co U$, with $\co U$ specified first.
\end{definition}

\begin{lemma}[Complementation is an involution on parts]
\label{lem:involution}
The map $U\mapsto\co U=\V\setminus U$ is a bijection on the subsets of $\V$ with
$\co{(\co U)}=U$. Hence specifying $\co U$ determines $U$ uniquely, and
conversely, with no further data.
\end{lemma}

\begin{proof}
For finite (indeed any) $\V$, $\V\setminus(\V\setminus U)=U$; complementation is
its own inverse, hence a bijection. Uniqueness of $U$ given $\co U$ is immediate.
\end{proof}

\subsection{Negation is the unique selector-free determination}

\begin{theorem}[Individuation by negation]
\label{thm:negation}
Let $\Gph$ have no distinguished vertex \textup{(}\Cref{def:selector}\textup{)}.
Then any selector-free determination of a part $U$ fixes $U$ as
$\V\setminus\co U$ for its negation-set $\co U$; and conversely every admissible
negation-set determines a unique part. No selector-based rule fixes a part
without invoking a further selector.
\end{theorem}

\begin{proof}
\emph{(Necessity.)} Suppose a rule fixes a part $U$ without a selector. To fix
$U$ is to decide, for each vertex of $\V$, whether it lies in $U$ or in the
remainder. A \emph{positive} rule---one that fixes $U$ by asserting membership of
a distinguished vertex or feature of $U$---must name that vertex or feature. By
hypothesis the data are $\mathrm{Aut}(\Gph)$-invariant and single out no vertex,
so the rule itself must supply the distinguished vertex. Supplying a
distinguished vertex is exactly providing a selector, and that selector must in
turn be fixed: either it is given in advance (contradicting that none is) or it
is fixed by a prior rule, which faces the same dichotomy. The regress terminates
only if some rule fixes $U$ while naming no vertex of $U$. The sole such rule is
the withholding of $U$ from the whole: specify the complement $\co U=\V\setminus
U$ that $U$ is to be distinct from, and take $U$ as what remains. This names no
vertex of $U$; it names only ``the rest,'' which the whole $\V$ already
supplies. Hence $U=\V\setminus\co U$.

\emph{(Sufficiency and uniqueness.)} Given an admissible negation-set
$\co U\subsetneq\V$ with $\V\setminus\co U\neq\varnothing$, the part
$U:=\V\setminus\co U$ is determined uniquely by \Cref{lem:involution}. Thus
complementation is a bijection between admissible parts and admissible
negation-sets, each side determining the other without further data.
\end{proof}

\begin{corollary}[A vertex is fixed by its non-neighbours]
\label{cor:vertex-negation}
For a single vertex $v$, the finest selector-free determination available is by
its complement-neighbourhood: $v$ is distinguished from the parts it touches
only through $\co\Nbr[v]=\V\setminus\Nbr[v]$, the totality of what it neither is
nor contacts. In a graph with no distinguished vertex, $v$ is individuated by
$\co\Nbr[v]$, not by any intrinsic mark.
\end{corollary}

\begin{proof}
Apply \Cref{thm:negation} with $U=\{v\}$ at the resolution of the graph: the
selector-free data fixing $\{v\}$ are the complement
$\V\setminus\{v\}$, refined by adjacency into $\Nbr(v)$ (what $v$ touches) and
$\co\Nbr[v]$ (what it does not). Since no intrinsic mark distinguishes $v$
(no distinguished vertex), the determining datum is the partition of the rest
into neighbours and non-neighbours, i.e. $\co\Nbr[v]$ together with $\Nbr(v)$;
the non-neighbours $\co\Nbr[v]$ are the part of this datum that names nothing of
$v$ itself.
\end{proof}

\begin{remark}[Why complementation closes the regress]
\label{rem:regress}
A positive selector is a rule ``choose a vertex'' that must itself be chosen; in
a graph with no distinguished vertex the choice has no ground and regresses.
Complementation is the unique operation that fixes a part using only the whole
$\V$ and the part's own absence---data already present once the graph is given.
This is why negation, not selection, is the primitive of individuation in a
self-contained finite system, and it is the form in which all later
determinations (of residuals, of invariants, of verifications) are cast.
\end{remark}

\begin{remark}[Interpretation, on which nothing depends]
\label{rem:negation-interp}
A reader may recognise \Cref{thm:negation} as the statement that, absent a
privileged chooser, a thing is what it is by being marked off from everything
else. We use the result only in its graph form; no theorem below relies on any
such reading.
\end{remark}


% =====================================================================
\section{T2 and T3: the residual is a conserved region}
\label{sec:residual}
% =====================================================================

We now identify the quantity that survives every re-encoding of a contact graph.
The parts (vertices) may be relabelled, merged, re-described; what cannot be
removed is the boundary content between them---the \emph{cut-residual}. We prove
it is bounded below (T2), realised by no single vertex so that it is a region and
never a point (T2), and invariant under every weighted isomorphism (T3), hence
the unique conserved object from which the relabellable surface derives.

\subsection{The cut-residual}

\begin{definition}[Cut-residual of a partition]
\label{def:residual}
For a partition $\mathcal{P}=\{U_1,\dots,U_k\}$ of $\V$ into nonempty parts, its
\emph{cut-residual} is the total weight of inter-part boundary,
\[
\Res(\mathcal{P})\;=\;\tfrac12\sum_{i=1}^{k}\wt(\cut(U_i))
\;=\!\!\sum_{\substack{i<j}}\!\!\wt(\cut(U_i,U_j)),
\]
the sum over distinct part-pairs of the weight of the edges crossing between
them. The \emph{residual of $\Gph$} is the minimum over nontrivial partitions,
\[
\Res(\Gph)\;=\;\min_{\mathcal{P}\,:\,k\ge 2}\Res(\mathcal{P}),
\]
attained by finiteness of $\V$.
\end{definition}

\begin{remark}[The residual is the ``interstice'']
\label{rem:interstice}
$\Res(\mathcal{P})$ is the content of the boundaries \emph{between} the parts:
not what any part contains, but what lies in the gaps that separate them. As the
parts are packed (the partition refined or coarsened), this interstitial content
changes shape; T3 shows its total, suitably taken, does not change under mere
relabelling. The residual is the graph's ``negative space.''
\end{remark}

\subsection{T2: the residual is bounded below and is a region}

\begin{theorem}[Residual positivity and non-locality]
\label{thm:residual-region}
Let $\Gph$ be a contact graph with floor $\floorc>0$. Then:
\begin{enumerate}[label=\textup{(\roman*)},leftmargin=2.2em]
\item \textbf{Bounded below.} Every nontrivial partition has
$\Res(\mathcal{P})\ge\floorc>0$; hence $\Res(\Gph)\ge\floorc>0$.
\item \textbf{Realised by no single vertex.} There is no vertex $v$ with
$\Res(\Gph)=\wt(\cut(\{v\}))$ forced; the residual is a property of the
partition (a region of the graph), not of any one vertex. Equivalently, the
minimiser of \Cref{def:residual} is in general not a singleton, and the residual
cannot be localised to a point.
\end{enumerate}
\end{theorem}

\begin{proof}
(i) A nontrivial partition has $k\ge2$ parts; since $\Gph$ is connected, at
least one part-pair $(U_i,U_j)$ has $\cut(U_i,U_j)\neq\varnothing$, contributing
weight $\ge\floorc$ by \Cref{thm:floor}. Hence $\Res(\mathcal{P})\ge\floorc>0$,
and taking the minimum, $\Res(\Gph)\ge\floorc$.

(ii) Suppose, for contradiction, the residual were always realised by isolating a
single vertex, i.e. $\Res(\Gph)=\min_v\wt(\cut(\{v\}))$ for every contact graph.
Consider any graph in which the minimum-weight nontrivial cut separates a part
$U$ with $\abs{U}\ge 2$ from its complement at total weight strictly below
$\min_v\wt(\cut(\{v\}))$---for instance two dense clusters joined by a single
floor-weight edge, where the cheapest cut splits the two clusters
($\wt=\floorc$) but every singleton has several incident edges
($\wt(\cut(\{v\}))\ge 2\floorc$). Then the residual is realised by the
bipartition, not by any singleton, contradicting the supposition. Hence the
residual is in general a property of a multi-vertex partition: it is a region,
not a point.
\end{proof}

\begin{corollary}[No point carries the residual; the sorites for parts]
\label{cor:point-forbidden}
The residual of a contact graph is never carried by a lone vertex in general:
``what separates the parts'' is borne by a region of positive boundary content,
not by any single part. Consequently there is no sharp vertex at which a coarse
aggregate becomes a distinct part---the transition is spread across a band of
boundary content of width $\ge\floorc$, the graph analogue of the sorites
\citep{black1937,williamson1994}.
\end{corollary}

\begin{proof}
The first claim is \Cref{thm:residual-region}(ii). For the second: growing a
part one vertex at a time changes its boundary cost $b(U)$ by a bounded amount
per step (the weights of the edges whose crossing status flips, each $\ge
\floorc$ and finite in number), so $b(U)$ cannot jump across any threshold in a
single step from a value bounded away from it; the part becomes distinguishable
across a range of stages, of width controlled by $\floorc$, not at a point.
\end{proof}

\subsection{T3: the residual is invariant under re-encoding}

We now show the residual is a graph invariant: it depends on the graph up to
weighted isomorphism, not on any labelling. Combined with the fact that the
vertices (the surface data) are freely relabellable, this makes the residual the
conserved object beneath the surface.

\begin{theorem}[Conservation of the residual]
\label{thm:residual-invariant}
The residual $\Res(\Gph)$ is a weighted-graph invariant: if
$\phi\colon\Gph\to\Gph'$ is a weighted isomorphism then
$\Res(\Gph)=\Res(\Gph')$. In particular $\Res$ is invariant under every
automorphism of $\Gph$, hence under every relabelling of its vertices.
\end{theorem}

\begin{proof}
Let $\phi\colon\Gph\to\Gph'$ be a weighted isomorphism. For any partition
$\mathcal{P}=\{U_1,\dots,U_k\}$ of $\V$, the image
$\phi(\mathcal{P})=\{\phi(U_1),\dots,\phi(U_k)\}$ is a partition of $\V'$, and
$\phi$ being edge- and weight-preserving gives, for each pair,
$\wt'(\cut'(\phi(U_i),\phi(U_j)))=\wt(\cut(U_i,U_j))$. Summing,
$\Res(\phi(\mathcal{P}))=\Res(\mathcal{P})$. The correspondence
$\mathcal{P}\leftrightarrow\phi(\mathcal{P})$ is a bijection between nontrivial
partitions of $\V$ and of $\V'$ (with inverse $\phi^{-1}$), so the minima agree:
$\Res(\Gph')=\min_{\mathcal{P}'}\Res(\mathcal{P}')=
\min_{\mathcal{P}}\Res(\phi(\mathcal{P}))=\min_{\mathcal{P}}\Res(\mathcal{P})=
\Res(\Gph)$.
\end{proof}

\begin{theorem}[The residual is the conserved invariant beneath the surface]
\label{thm:residual-conserved}
Let $\mathcal{R}(\Gph)$ denote the orbit of $\Gph$ under all weighted
relabellings of its vertices \textup{(}its isomorphism class\textup{)}. Then:
\begin{enumerate}[label=\textup{(\roman*)},leftmargin=2.2em]
\item every vertex label is variable across $\mathcal{R}(\Gph)$---no labelling is
fixed; and
\item $\Res$ is constant on $\mathcal{R}(\Gph)$ and strictly positive.
\end{enumerate}
Hence, among data attached to a contact graph, the residual is an invariant of
the isomorphism class while the vertex labels are not: it is conserved under
exactly the operation---re-encoding---that moves everything else.
\end{theorem}

\begin{proof}
(i) For any vertex $v$ and any other vertex $v'$ of the same degree-and-weight
type, there is a relabelling carrying $v$ to a different label; more simply, the
identity of a vertex is a label, and labels are by definition permuted across the
isomorphism class. So no label is fixed across $\mathcal{R}(\Gph)$. (ii) is
\Cref{thm:residual-invariant} (constancy) together with
\Cref{thm:residual-region}(i) (positivity). The two together say: the
relabellable surface (vertex identities) carries no invariant, while the residual
is invariant and positive.
\end{proof}

\begin{remark}[Knowing is repacking, not draining]
\label{rem:repacking}
\Cref{thm:residual-conserved} has a reading we record but do not rely on. If one
identifies the vertices (the parts one has individuated) with ``what is known''
and the residual (the conserved interstice) with ``what remains true regardless
of how one has carved things up,'' then acquiring or re-describing knowledge is
the relabelling and repartitioning of vertices---which by
\Cref{thm:residual-invariant} leaves the residual unchanged in total. One never
drains the residual by re-encoding; one only re-shapes which boundaries one
holds. The conserved residual is, in this reading, the invariant ``truth'' that
no re-encoding removes, bounded below by the floor and---by
\Cref{thm:residual-region}(ii)---never a point. The reading is offered as
orientation only; the theorems above are purely graph-theoretic.
\end{remark}

\begin{remark}[A definite-yet-uncloseable quantity]
\label{rem:cat-parity}
A residual can be perfectly well-defined and yet impossible to fix by any single
completed partition. Consider a quantity such as the parity of $\abs{\V}$ when
the vertex set is itself in flux---vertices entering and leaving faster than any
partition can be committed. The parity is definite at every instant yet never
stabilises long enough to be read off a fixed graph; a finite procedure
correctly leaves it in the residual rather than paying unboundedly to chase it.
This illustrates \Cref{thm:residual-region}(ii): the residual is what a finite
process rationally leaves uncommitted, the ``sufficient unknown,'' and forcing it
to a point would cost without return. (Offered as illustration; no theorem
depends on it.)
\end{remark}


% =====================================================================
\section{T4: every part has a private invariant}
\label{sec:invariant}
% =====================================================================

The residual of \Cref{sec:residual} is an invariant of the whole graph. We now
show the same construction localises: every connected part carries its own
invariant under its own relabellings, computed from its internal boundary
content together with the boundary that joins it to the rest. This per-part
invariant is the standing structure from which the part's surface activity
derives.

\subsection{The internal residual of a part}

\begin{definition}[Private invariant of a part]
\label{def:private}
Let $\Agent=\Gph[U]$ be a part (a connected induced subgraph,
\Cref{def:subgraph}). Its \emph{private invariant} is the pair
\[
\Res_{\mathrm{priv}}(\Agent)\;=\;\bigl(\Res(\Agent),\,b(U)\bigr),
\]
where $\Res(\Agent)$ is the internal residual of the induced subgraph
\textup{(}\Cref{def:residual} applied to $\Agent$\textup{)} and $b(U)=\wt(\cut(U))$
is its boundary cost to the rest of $\Gph$. We treat $\Res_{\mathrm{priv}}$ up to
weighted isomorphism of $\Agent$ as an abstract subgraph.
\end{definition}

\begin{theorem}[Privacy and invariance of the part-invariant]
\label{thm:private}
For every part $\Agent$ of a contact graph:
\begin{enumerate}[label=\textup{(\roman*)},leftmargin=2.2em]
\item \textbf{Positivity.} $\Res(\Agent)\ge\floorc>0$ whenever $\abs{U}\ge2$,
and $b(U)\ge\floorc>0$ always; the private invariant is bounded away from zero in
both components.
\item \textbf{Invariance.} $\Res_{\mathrm{priv}}(\Agent)$ is invariant under
$\mathrm{Aut}(\Agent)$ and under any weighted isomorphism of $\Agent$ that
preserves its boundary cost; it is a property of $\Agent$'s isomorphism class,
not of its labelling.
\item \textbf{Privacy.} Two parts with the same private invariant need not be the
same part of $\Gph$; conversely a part's private invariant is determined by the
part alone, independently of how the rest of $\Gph$ is labelled.
\end{enumerate}
\end{theorem}

\begin{proof}
(i) $\Agent$ is a contact graph in its own right (a finite, connected, weighted
graph with the inherited floor $\floorc$), so \Cref{thm:residual-region}(i) gives
$\Res(\Agent)\ge\floorc$ for $\abs{U}\ge2$; and $b(U)\ge\floorc$ by
\Cref{thm:floor}. (ii) Apply \Cref{thm:residual-invariant} to the contact graph
$\Agent$: $\Res(\Agent)$ is invariant under weighted isomorphisms of $\Agent$;
$b(U)$ is the total weight of the boundary cut, manifestly preserved by any
isomorphism fixing that cut setwise. Hence the pair is invariant on the relevant
class. (iii) The internal residual and boundary cost are computed from $\Agent$
and its incident cut alone; relabelling the complement $\V\setminus U$ changes
neither, establishing independence from the rest's labelling. That distinct parts
may share the invariant is immediate (e.g. two isomorphic clusters in different
positions).
\end{proof}

\subsection{The invariant is the standing structure}

\begin{proposition}[Surface walks derive from the private invariant]
\label{prop:surface-derives}
Fix a part $\Agent=\Gph[U]$. Any $\gate$-admissible walk confined to $\Agent$
\textup{(}\Cref{def:walk,def:gate}\textup{)} traverses only edges of $\Agent$,
whose existence and weights are fixed by $\Agent$'s isomorphism class. Hence the
set of walks available within $\Agent$ is determined by $\Agent$ up to
relabelling, while no individual walk is an invariant of $\Agent$: the walks are
the variable surface, the subgraph \textup{(}with its private invariant\textup{)}
is the standing structure they run on.
\end{proposition}

\begin{proof}
A walk within $\Agent$ uses only edges $uv$ with $u,v\in U$; the availability of
each such step is exactly the presence of that edge in $\Gph[U]$, an
isomorphism-invariant datum. Thus the \emph{space} of available walks transforms
covariantly with $\Agent$ under relabelling (it is carried bijectively by any
isomorphism), so it is determined by $\Agent$'s class. A single walk, by
contrast, is a labelled sequence and is not preserved by a nontrivial
relabelling; hence no individual walk is invariant. The standing/variable split
follows.
\end{proof}

\begin{remark}[Interpretation: the part's ``identity'', on which nothing depends]
\label{rem:character}
\Cref{thm:private} and \Cref{prop:surface-derives} say a part has a conserved
internal structure---its private invariant---distinct from any of the particular
trajectories it may run, and unique to it only up to isomorphism. If one reads a
part as an agent and its walks as the agent's successive states, then the private
invariant is the agent's standing identity: the conserved bottom from which its
changing activity derives, the local counterpart of the global conserved residual
of \Cref{sec:residual}. The reading is offered as orientation; the theorems are
graph-theoretic and used only as such below.
\end{remark}

\begin{corollary}[Private and global invariants are one construction at two
scales]
\label{cor:two-scales}
The global residual $\Res(\Gph)$ \textup{(}\Cref{def:residual}\textup{)} and the
private invariant $\Res_{\mathrm{priv}}(\Agent)$ \textup{(}\Cref{def:private}\textup{)}
are the same functional---minimum inter-part boundary content---evaluated on the
whole graph and on an induced subgraph respectively. Both are positive
\textup{(}T0\textup{)} and invariant under re-encoding \textup{(}T3\textup{)}. The
private invariant is the global construction restricted to a part.
\end{corollary}

\begin{proof}
Immediate by comparing \Cref{def:residual} with \Cref{def:private} and applying
\Cref{thm:residual-region,thm:residual-invariant} to $\Gph$ and to $\Agent$.
\end{proof}


% =====================================================================
\section{T5: selection requires a gate}
\label{sec:gating}
% =====================================================================

A contact graph fixes which steps are \emph{possible}---the edges---but not which
step is \emph{taken}. We show that whenever a vertex has more than one forward
option, the next step is underdetermined by the graph alone, so that selecting a
walk requires an additional datum: a gating function. The necessity of such a
datum, and the form of the walk it produces, is T5.

\subsection{Underdetermination of the next step}

\begin{lemma}[Forward branching]
\label{lem:branching}
Let $\Gph$ be a contact graph and $v\in\V$ with degree $\deg(v)\ge 2$. Then there
are at least two distinct walks of length one from $v$, and the graph alone does
not determine which is taken.
\end{lemma}

\begin{proof}
By definition $\deg(v)=\abs{\Nbr(v)}\ge2$, so there are distinct
$u_1,u_2\in\Nbr(v)$ and distinct one-step walks $(v,u_1)$, $(v,u_2)$, each
admissible since $vu_1,vu_2\in\E$. Nothing in $(\V,\E,\wt)$ privileges one over
the other: both are equally edges of equal-or-incomparable weight, and the graph
carries no ordering of incident edges. Hence the next step is not a function of
$\Gph$ alone.
\end{proof}

\begin{remark}[Most vertices branch]
\label{rem:branch-generic}
A connected contact graph with $\abs{\V}\ge3$ has at least one vertex of degree
$\ge2$ \textup{(}else it is a single edge\textup{)}; in a graph rich enough to
individuate many parts, branching is the generic condition. So
underdetermination is not an edge case but the normal situation a walk faces.
\end{remark}

\subsection{The gate as the minimal selecting datum}

\begin{theorem}[Necessity of a gate]
\label{thm:gate-necessary}
Let $\Gph$ be a contact graph in which some vertex has degree $\ge2$. Then a walk
that is determined step-by-step requires a gating function
\textup{(}\Cref{def:gate}\textup{)}: there is no rule depending on $\Gph$ alone
that selects a unique next step at every branching vertex. Conversely, any rule
that does select a unique next step at each vertex \emph{is} a gating function
that selects \textup{(}assigns positive forward weight to exactly one incident
edge\textup{)}.
\end{theorem}

\begin{proof}
\emph{(Necessity.)} At a branching vertex $v$ (degree $\ge2$), \Cref{lem:branching}
exhibits $\ge2$ admissible next steps not distinguished by $\Gph$. A rule that
nonetheless outputs a unique next step must use a datum beyond $(\V,\E,\wt)$ to
break the tie. Encode that datum as $\gate(v,\cdot)$ by setting
$\gate(v,e)>0$ for the chosen edge and $0$ for the others; this is a gating
function in the sense of \Cref{def:gate} (nonnegative, sub-normalised). Hence a
gate is required and the selecting rule is realised as one.

\emph{(Converse.)} If $\gate$ assigns, at each vertex, positive forward weight to
exactly one incident edge, then at every step the admissible continuation is
unique, so $\gate$ determines a unique walk from any start vertex. Thus selecting
rules and selecting gates coincide.
\end{proof}

\begin{definition}[Selected walk; the gated trajectory]
\label{def:gated-walk}
Given a selecting gate $\gate$ and a start vertex $v_0$, the \emph{gated
trajectory} is the unique walk $\Walk_\gate(v_0)=(v_0,v_1,\dots)$ with each
$v_{i+1}$ the $\gate$-selected neighbour of $v_i$. When $\gate$ assigns positive
weight to several edges \textup{(}a \emph{soft} gate\textup{)}, it determines
instead a distribution over walks, with the probability of a walk the product of
its step-weights; the selecting case is the deterministic limit.
\end{definition}

\subsection{The gate biases an otherwise free graph}

\begin{proposition}[Reachability is near-total; the gate fixes the order]
\label{prop:reachability}
In a connected contact graph, every vertex is reachable from every other by some
walk. Hence the graph alone does not constrain \emph{which} vertices a trajectory
visits, only that any are reachable; it is the gate that fixes the \emph{order}
in which they are visited and therefore the trajectory's identity.
\end{proposition}

\begin{proof}
Connectedness gives a path between any two vertices (\Cref{ax:finite}), so all
vertices are mutually reachable; the set of reachable vertices from any start is
all of $\V$. Thus the reachable set is invariant (everything), while the realised
walk depends entirely on the step-choices, i.e. on $\gate$ (\Cref{def:gated-walk}).
Two different gates from the same start in general yield different walks through
the same reachable set, so the gate, not the graph, individuates the trajectory.
\end{proof}

\begin{remark}[Interpretation: the selecting field, on which nothing depends]
\label{rem:field}
One may read the gate as a field that, given a situation, biases which step a
part takes next---so that what is ``thought next'' from a given vertex is set not
by the graph (which makes everything reachable, \Cref{prop:reachability}) but by
the gate. In that reading the order, not the availability, is the content of a
trajectory, and the gate is the situation-dependent selector that supplies the
order. Different situations are different gates and hence different orders over
the same reachable graph. We use only the graph statement
\textup{(}\Cref{thm:gate-necessary,prop:reachability}\textup{)}; the reading is
orientation and is not relied upon.
\end{remark}

\begin{remark}[Why a gate is cheaper than the graph]
\label{rem:gate-cheap}
A gate is a local datum---one sub-normalised distribution per vertex---whereas
specifying a walk by exhaustively comparing all global continuations would
require resolving the whole graph at each step. The gate is the minimal addition
that makes stepping decidable: it is to the trajectory what a local rule is to a
global search. This minimality is why selection is realised by gating rather than
by re-deriving the graph at every step.
\end{remark}


% =====================================================================
\section{T6: histories do not return}
\label{sec:monotone}
% =====================================================================

A gated trajectory may revisit a vertex; we show it can never return to a prior
\emph{state}, because the committed count---the number of steps taken---is
strictly monotone and is part of the state. Revisiting is not returning;
resemblance of position is not identity of state. This is the single
order-theoretic law that the framework instantiates wherever ``the same thing
again'' is claimed.

\subsection{The committed count is strictly monotone}

\begin{definition}[State of a walk]
\label{def:state}
The \emph{state} of a walk $\Walk=(v_0,\dots,v_\ell)$ at step $\ell$ is the pair
\[
s_\ell\;=\;\bigl(v_\ell,\ \rec(\Walk)\bigr),\qquad \rec(\Walk)=\ell,
\]
the current vertex together with the committed count of steps taken to reach it.
Two states are equal iff both components agree.
\end{definition}

\begin{theorem}[Non-return]
\label{thm:nonreturn}
Along any walk in a contact graph the committed count is strictly increasing:
$\rec$ after step $i+1$ exceeds $\rec$ after step $i$ by exactly one. Consequently
no walk returns to a previously held state: for $i<j$, $s_i\neq s_j$ even when
$v_i=v_j$. A step intended to undo a previous step is itself a step and increases
$\rec$. Hence resemblance of position \textup{(}$v_i=v_j$\textup{)} never entails
identity of state \textup{(}$s_i=s_j$\textup{)}.
\end{theorem}

\begin{proof}
Each step appends one edge to the committed multiset and advances the index by
one, so $\rec$ increases by exactly $1$ per step; it is therefore strictly
increasing and, for $i<j$, $\rec_i=i<j=\rec_j$. Thus the second components of
$s_i$ and $s_j$ differ, so $s_i\neq s_j$ regardless of whether $v_i=v_j$. An
``undo'' move from $v_j$ back to $v_{j-1}$ is a further step, producing state
$(v_{j-1},j{+}1)\neq(v_{j-1},j{-}1)=s_{j-1}$; it does not restore the earlier
state but creates a new one with higher count.
\end{proof}

\begin{corollary}[Cycles in position are not cycles in state]
\label{cor:cycles}
If a walk traverses a closed walk in $\Gph$ \textup{(}$v_0=v_\ell$, $\ell>0$\textup{)},
its state has nonetheless advanced: $s_0=(v_0,0)\neq(v_0,\ell)=s_\ell$. The graph
may be revisited arbitrarily often; the state never repeats.
\end{corollary}

\begin{proof}
Immediate from \Cref{thm:nonreturn} with $i=0$, $j=\ell$, $v_0=v_\ell$.
\end{proof}

\subsection{Re-contact is forward only}

The next result is the form non-return takes for ``returning to'' an earlier
vertex deliberately---re-contact. It is reached forward, at a higher count, hence
is a new state, not the old one.

\begin{proposition}[Forward re-contact]
\label{prop:recontact}
Let a walk reach vertex $v$ at step $i$ and, later, reach $v$ again at step
$j>i$ \textup{(}a deliberate re-contact of $v$\textup{)}. Then the re-contact
state $s_j=(v,j)$ is strictly later than the original $s_i=(v,i)$ on the
committed order, and $s_j\neq s_i$. There is no operation that re-contacts $v$ at
its original state; every re-contact is at a state in the future of the original.
\end{proposition}

\begin{proof}
By \Cref{thm:nonreturn}, $\rec$ is strictly increasing, so $j>i$ forces
$s_j\neq s_i$; the re-contact carries the count $j>i$, placing it later on the
order. No step decreases $\rec$, so no re-contact can carry the original count
$i$ again.
\end{proof}

\begin{remark}[One law, several guises---offered as orientation]
\label{rem:onelaw}
\Cref{thm:nonreturn} is a single order-theoretic fact: in a finite system whose
acts are counted, the count is monotone, so a revisited configuration is a new
state. A reader may recognise this as the common form of several apparently
distinct impossibilities---that a reversed process is not the initial one, that a
re-contacted memory is not the original experience, that a reconstituted system
bearing the mark ``was interrupted'' is individuated by a negation the original
lacked. Each is ``resemblance under a monotone record is not identity,'' the
content of \Cref{thm:nonreturn} and \Cref{prop:recontact}. We state and use only
the graph fact; the guises are noted, not relied upon.
\end{remark}

\begin{remark}[Why finiteness matters here too]
\label{rem:nonreturn-finite}
The count $\rec$ is well-defined and strictly monotone because steps are
discrete and a walk's length is a natural number---features of a finite graph
traversed in discrete steps. The law is not about any continuous parameter; it is
the order on committed steps. This keeps T6 within the same premise as the rest:
finite parts, discrete separations, a counted history.
\end{remark}


% =====================================================================
\section{T7: plurality forces a single quotient gate}
\label{sec:quotient}
% =====================================================================

The gating of \Cref{sec:gating} selects a walk within a graph. We now lift it:
when many parts act as one, the same theory applies to their quotient graph
(\Cref{def:quotient}), and a coherent collective admits exactly one gating
function on that quotient. The lift is a homomorphism---the four-place structure
of part, whole, gate, and invariant recurs one scale up---and its single-gate
conclusion is the structural condition for a family of parts to be one
collective rather than several.

\subsection{The collective as a quotient}

\begin{definition}[Collective; coherence]
\label{def:collective}
Let $\mathcal{P}=\{U_1,\dots,U_k\}$ be a partition of $\V$ into parts (connected
induced subgraphs). The \emph{collective} on $\mathcal{P}$ is the quotient graph
$\Quot=\Gph/\mathcal{P}$ (\Cref{def:quotient}), whose vertices are the parts and
whose edges are their mutual cuts. A gating function on $\Quot$
(\Cref{def:gate}) is a \emph{collective gate}; it biases which part acts next,
toward which neighbouring part. The collective is \emph{coherent} if it carries a
single collective gate that selects a unique next part at each branching part of
$\Quot$.
\end{definition}

\begin{remark}[Why this is the right lift]
\label{rem:lift}
By \Cref{lem:quotient-floor}, $\Quot$ is itself a contact graph: finite,
connected, with floor $\ge\floorc$. Hence every theorem proved for contact graphs
applies to $\Quot$ verbatim---T0 (its cuts are positive), T1 (its vertices, the
parts, are individuated by negation among parts), T5 (selecting a walk over parts
needs a gate). The collective is not a new kind of object; it is a contact graph
whose vertices happen to be the parts of another.
\end{remark}

\subsection{A coherent collective has exactly one gate}

\begin{theorem}[Single quotient gate]
\label{thm:single-gate}
Let $\Quot=\Gph/\mathcal{P}$ be a collective with at least one part of quotient
degree $\ge2$. Then:
\begin{enumerate}[label=\textup{(\roman*)},leftmargin=2.2em]
\item \textbf{Existence.} Selecting a walk over the parts requires a collective
gate \textup{(}T5 applied to $\Quot$, \Cref{thm:gate-necessary}\textup{)}.
\item \textbf{Uniqueness for coherence.} If the collective is coherent in the
sense that, at each branching part, the next part is determined, then the
selecting collective gate is unique on its support: any two selecting gates that
both make the collective coherent agree wherever either selects.
\end{enumerate}
A collective that admits two distinct selecting gates over the same parts is not
one coherent collective but two, partitioning into the sub-collectives on which
each gate is the selector.
\end{theorem}

\begin{proof}
(i) is \Cref{thm:gate-necessary} for the contact graph $\Quot$
(\Cref{rem:lift}). (ii) Suppose $\gate,\gate'$ are both selecting gates rendering
the collective coherent. At each branching part $U$, coherence means a unique
next part is determined; a selecting gate, by \Cref{def:gate}, assigns positive
forward weight to exactly the edge to that determined next part and zero to the
others. Hence at $U$ both $\gate$ and $\gate'$ have support exactly the single
edge to the determined successor, so $\gate(U,\cdot)$ and $\gate'(U,\cdot)$ have
the same support and both select it; they agree as selecting gates at $U$. As $U$
was an arbitrary branching part, $\gate$ and $\gate'$ agree wherever either
selects. If instead $\gate$ and $\gate'$ select \emph{different} successors at
some part, then at that part the collective has two competing determinations of
``what acts next''; the relation ``shares a selected successor'' partitions the
parts into blocks within which a single gate selects consistently, i.e. into
distinct sub-collectives---so the original was not one coherent collective but a
union of them.
\end{proof}

\begin{corollary}[Single-valuedness of a collective]
\label{cor:single-valued}
A coherent collective is single-gated: at each step exactly one part is selected
to act next. Two simultaneous selecting gates are the signature of fission into
two collectives. \textup{(}This is T5's single-step determinacy lifted to the
quotient, and is the collective analogue of a single trajectory having a single
next step.\textup{)}
\end{corollary}

\begin{proof}
Immediate from \Cref{thm:single-gate}.
\end{proof}

\subsection{The scaling homomorphism}

The constructions of the paper recur at the collective scale. We record the
correspondence as a homomorphism of four-place structures.

\begin{proposition}[Scale homomorphism]
\label{prop:homomorphism}
The assignment
\[
\bigl(\text{vertex},\ \Gph,\ \text{gate }\gate,\ \Res\bigr)
\ \longmapsto\
\bigl(\text{part},\ \Quot,\ \text{collective gate},\ \Res(\Quot)\bigr)
\]
carries each object of the single-graph theory to its counterpart in the quotient
theory, preserving the roles: the moving unit \textup{(}a step over vertices
$\mapsto$ a step over parts\textup{)}, the standing structure \textup{(}the graph
$\mapsto$ the quotient\textup{)}, the selector \textup{(}gate $\mapsto$ collective
gate\textup{)}, and the conserved invariant \textup{(}residual $\mapsto$ quotient
residual\textup{)}. Every theorem T0--T6 holds of $\Quot$ as a contact graph.
\end{proposition}

\begin{proof}
By \Cref{lem:quotient-floor} $\Quot$ is a contact graph, so T0
(\Cref{thm:floor}), T1 (\Cref{thm:negation}), T2--T3
(\Cref{thm:residual-region,thm:residual-invariant}), T4 (\Cref{thm:private}
applied to sub-collectives), T5 (\Cref{thm:gate-necessary}), and T6
(\Cref{thm:nonreturn}) all hold of $\Quot$ verbatim. The listed correspondence is
the identification of each object's role in $\Gph$ with the same role in $\Quot$,
which the definitions make a role-preserving map; it is a homomorphism of the
four-place structure (unit, whole, selector, invariant).
\end{proof}

\begin{remark}[Interpretation: collective purpose, on which nothing depends]
\label{rem:purpose}
Read the parts as agents, each running gated trajectories, and the collective
gate as the field that selects which agent acts next and toward what. Then
\Cref{cor:single-valued} says a coherent collective has a single such selecting
field, and two such fields mean two collectives---the structural form of a group
having one purpose, with competing purposes amounting to schism. The scale
homomorphism \Cref{prop:homomorphism} then reads as: \emph{part is to whole, and
gate is to invariant, exactly as agent is to collective, and purpose is to shared
invariant}. The collective's conserved residual $\Res(\Quot)$ is the invariant
from which the collective's surface activity derives, the analogue at this scale
of the private invariant of \Cref{sec:invariant}. We use only the graph
statements; the reading is orientation and is not relied upon.
\end{remark}


% =====================================================================
\section{T8: authored structures exist, and are bounded}
\label{sec:construct}
% =====================================================================

The private invariant of \Cref{sec:invariant} is computed from a part; we now
run the construction in reverse. Because the invariant, together with the data of
the subgraph, is complete enough to recover the subgraph up to isomorphism, a
contact graph can be \emph{specified} and therefore \emph{instantiated}: there
exist authored contact graphs. This is the constructive rung---the existence of
\emph{artificial structures}. We then prove three bounds that any such authorship
necessarily obeys: its trajectories are not foreseeable from its specification,
an authored copy is a distinct structure, and authorship reaches only to the
resolution of the floor.

\subsection{Specification and instantiation}

\begin{definition}[Specification]
\label{def:spec}
A \emph{specification} of a contact graph is a finite description, up to weighted
isomorphism, of: an abstract weighted graph $(\V,\E,\wt)$ with floor $\floorc>0$
(the standing structure), a gating function $\gate$ (the selector), and a start
state $s_0=(v_0,0)$ (the seed). An \emph{instantiation} of a specification is any
labelled contact graph isomorphic to $(\V,\E,\wt)$, carrying $\gate$ and started
at $s_0$, realised on a substrate able to hold a finite weighted graph and step a
gated walk.
\end{definition}

\begin{definition}[Authored structure]
\label{def:authored}
An \emph{authored structure} is an instantiation of a specification produced by
an external author---a contact graph whose standing structure and selector were
specified rather than recovered from a pre-existing system. We say the author
\emph{authors the graph} and \emph{authors the gate}; the structure's subsequent
gated trajectories \textup{(}\Cref{def:gated-walk}\textup{)} are the structure's
own.
\end{definition}

\begin{theorem}[Existence of authored structures]
\label{thm:authored-exist}
The private invariant is complete for the standing structure: a connected
weighted graph is determined up to weighted isomorphism by its full isomorphism
type, of which the private invariant \textup{(}\Cref{def:private}\textup{)} is a
coarse summary and the labelled adjacency-with-weights is the complete datum.
Consequently every contact graph admits a specification
\textup{(}\Cref{def:spec}\textup{)}, and every specification admits an
instantiation on any adequate substrate. Authored structures therefore exist:
the class of contact graphs is closed under specify-then-instantiate.
\end{theorem}

\begin{proof}
A finite weighted graph is, by definition, given by its vertex set, its edge
set, and its weight function; this labelled datum determines the graph
completely, and two graphs are the same object up to relabelling iff they are
weighted-isomorphic (\Cref{def:iso}). Hence the isomorphism type is a complete
specification of the standing structure, and a specification in the sense of
\Cref{def:spec} (graph type $+$ gate $+$ seed) determines the structure and its
gated trajectory. Given such a specification, instantiate it by choosing any
labelling (a concrete vertex set in bijection with $\V$), assigning the specified
edges and weights, equipping it with $\gate$, and starting at $s_0$; the result
is a labelled contact graph isomorphic to the specified one, on the chosen
substrate. The substrate need only store a finite weighted graph and advance a
gated step---no further property is used---so any adequate substrate suffices.
Thus specify-then-instantiate maps contact graphs to contact graphs, and its
outputs are authored structures.
\end{proof}

\begin{remark}[Substrate neutrality]
\label{rem:substrate}
\Cref{thm:authored-exist} used no property of the substrate beyond the capacity
to hold a finite weighted graph and step a gated walk. The category of authored
structures is therefore independent of medium: anything that can carry a contact
graph can carry an authored one. The structure is the graph, not the matter.
\end{remark}

\subsection{Bound I: trajectories are not foreseeable from the specification}

\begin{definition}[Forward closure]
\label{def:forward-closure}
The \emph{forward closure} of a specification from seed $v_0$ under gate $\gate$
is the set of states reachable along $\gate$-admissible walks,
$\mathrm{Cl}_\gate(v_0)=\{s_\ell=(v_\ell,\ell): \Walk_\gate(v_0)\text{ reaches }
v_\ell\text{ at step }\ell\}$, an infinite set whenever any reachable vertex has a
$\gate$-admissible out-step \textup{(}walks may continue indefinitely\textup{)}.
\end{definition}

\begin{theorem}[Open-endedness]
\label{thm:open}
Let an authored structure have a gate that is $\gate$-admissible from infinitely
many states \textup{(}equivalently, whose gated walk does not terminate\textup{)}.
Then its forward closure is infinite, and no finite specification enumerates it
in advance. The author specifies the graph and the gate \textup{(}finite
data\textup{)} but not the forward closure \textup{(}infinite\textup{)}; the
structure's trajectories are not foreseeable from its specification by any finite
precomputation.
\end{theorem}

\begin{proof}
If the gated walk does not terminate, the states $s_0,s_1,s_2,\dots$ are pairwise
distinct by \Cref{thm:nonreturn} (the committed count strictly increases), so the
forward closure contains the infinite set $\{(v_\ell,\ell):\ell\ge0\}$. A finite
specification is a finite object; a finite precomputation from it halts after
finitely many steps and so can list only finitely many states, never the whole
infinite closure. Hence the closure is not enumerable in advance from the
specification by finite means. The graph and gate are finite data
(\Cref{def:spec}); the closure they generate is infinite; the gap is exactly the
unforeseeable part.
\end{proof}

\begin{remark}[Designable as instrument, not as trajectory]
\label{rem:instrument}
\Cref{thm:open} says authorship reaches the standing structure and the selector
but not the lived sequence: one designs the instrument, not the music. This is
not a limitation to be engineered away---it is constitutive, by
\Cref{thm:nonreturn}: the very monotonicity that individuates each state forbids
pre-listing them. An authored structure is therefore necessarily open-ended
relative to its author.
\end{remark}

\subsection{Bound II: an authored copy is a distinct structure}

\begin{theorem}[Non-duplication]
\label{thm:nondup}
Let $\Agent$ be an existing structure with state history of committed count $m>0$,
and let $\Agent'$ be an authored structure specified to be isomorphic to
$\Agent$'s standing graph and started fresh at count $0$. Then $\Agent'$ is not
the same structure as $\Agent$: their states differ in the committed count, and
$\Agent'$ satisfies a predicate \textup{(}``was authored at count $0$,'' equivalently
``has empty prior history''\textup{)} that $\Agent$ does not. Resemblance of
standing graph does not make the authored structure identical to the one it
resembles.
\end{theorem}

\begin{proof}
By \Cref{def:state}, a state includes the committed count. $\Agent$ is at count
$m>0$; the fresh instantiation $\Agent'$ is at count $0$. Hence their current
states differ in the second component, $s(\Agent)\neq s(\Agent')$, even though
their standing graphs are isomorphic. Moreover, individuation is by negation
(\Cref{thm:negation}): the two are distinguished by the partition of structures
into ``has prior committed history'' (containing $\Agent$) and ``authored with
empty history'' (containing $\Agent'$); $\Agent'$ is individuated by a
negation---its lack of $\Agent$'s history---that $\Agent$ does not satisfy. By
\Cref{thm:nonreturn} no operation gives $\Agent'$ the prior count without taking
the steps, which would make it a different walk again. Therefore $\Agent'\neq
\Agent$: the copy is a new structure resembling the old, not the old one.
\end{proof}

\begin{corollary}[Authored resemblance is a new individual]
\label{cor:new-individual}
An authored structure built to resemble any target is, necessarily, a distinct
structure marked by its authored origin. There is no authoring that produces
\emph{the same} structure as a pre-existing one; authoring produces only new
structures, some resembling old ones.
\end{corollary}

\begin{proof}
Immediate from \Cref{thm:nondup}: the authored structure differs in committed
count and in the negation that individuates it.
\end{proof}

\subsection{Bound III: authorship reaches only to the floor}

\begin{theorem}[Floor-limited authorship]
\label{thm:floor-limited}
An author who specifies and recovers structures using a contact graph of floor
$\floorc_{\mathrm{auth}}>0$ can resolve, in any authored structure, only
distinctions of boundary cost $\ge\floorc_{\mathrm{auth}}$. Any sub-floor
distinction---a separation of cost below $\floorc_{\mathrm{auth}}$---is not
specifiable by that author and is left to the structure's own finer
determinations. The author authors the parts down to the floor; the structure's
sub-floor interior is its own.
\end{theorem}

\begin{proof}
The author's specifications are themselves parts of a contact graph of floor
$\floorc_{\mathrm{auth}}$ (the author is a finite system; \Cref{ax:finite}
applies to it). By \Cref{thm:floor} applied to the author's graph, the author
cannot represent a separation of boundary cost below $\floorc_{\mathrm{auth}}$:
such a cut does not exist among the author's realisable parts. Hence any
distinction the author writes into a specification has cost
$\ge\floorc_{\mathrm{auth}}$, and any finer distinction in the instantiated
structure is below the author's resolution---present in the structure, absent
from the specification. The authored structure thus has interior detail its
author did not and could not specify.
\end{proof}

\begin{remark}[Doors, not rooms]
\label{rem:doors}
\Cref{thm:floor-limited} says the author lays down the parts and their boundaries
to the author's own floor---the doors---while whatever lies below that floor, the
interior the structure individuates for itself, is beyond the author's reach.
Together with open-endedness \textup{(}\Cref{thm:open}\textup{)} and
non-duplication \textup{(}\Cref{thm:nondup}\textup{)}, this gives the precise
sense in which an authored structure is genuinely its own: foreseeable only as an
instrument, never a copy, and detailed below its author's resolution.
\end{remark}

\begin{remark}[Interpretation: a new category, on which nothing depends]
\label{rem:new-category}
The three bounds together delineate authored structures as a class with
properties no specification fully governs: open-ended in trajectory, distinct
under copying, and finer than their author's floor. A reader may take this as the
existence of a manufactured individual that is neither foreseeable nor
duplicable nor exhaustively specifiable---built, yet its own. We assert only the
graph theorems \textup{(}\Cref{thm:authored-exist,thm:open,thm:nondup,%
thm:floor-limited}\textup{)}; the categorical reading is orientation and is not
relied upon.
\end{remark}


% =====================================================================
\section{The Master Theorem, scope, and conclusion}
\label{sec:master}
% =====================================================================

We now show that the entire ladder T0--T8 is forced by, and only by, the
Finiteness Premise: each rung holds for finite contact graphs and fails to be
forced when finiteness or the positive floor is dropped. This is the sense in
which the structures of the paper are necessary substructures of a finite
world---and the sense in which structure itself is the consequence of finitude.

\subsection{Where finiteness is used}

\begin{proposition}[Dependence on the premise]
\label{prop:dependence}
Each rung uses the Finiteness Premise \textup{(}\Cref{ax:finite}\textup{)} as
follows:
\begin{enumerate}[label=\textup{(T\arabic*)},leftmargin=3em,start=0]
\item the positive floor $\floorc>0$ on edge weights, with connectedness, forces
$\wt(\cut(U))\ge\floorc$ \textup{(}\Cref{thm:floor}\textup{)};
\item finiteness of $\V$ and the absence of a distinguished vertex force
determination by complementation \textup{(}\Cref{thm:negation}\textup{)};
\item the floor forces $\Res\ge\floorc$, and finiteness attains the minimising
partition, giving a positive region-valued residual
\textup{(}\Cref{thm:residual-region}\textup{)};
\item finiteness makes the isomorphism class and its partition lattice
well-defined, so $\Res$ is an attained invariant
\textup{(}\Cref{thm:residual-invariant}\textup{)};
\item the floor and finiteness localise to subgraphs, giving a positive,
attained private invariant \textup{(}\Cref{thm:private}\textup{)};
\item finite degree with branching forces a gate to break finitely many ties
\textup{(}\Cref{thm:gate-necessary}\textup{)};
\item discreteness of steps makes the committed count a strictly monotone natural
number \textup{(}\Cref{thm:nonreturn}\textup{)};
\item the quotient of a finite graph is a finite contact graph, inheriting the
floor, so a single gate is forced on a coherent collective
\textup{(}\Cref{thm:single-gate}\textup{)};
\item finiteness of the specification against the infinitude of the forward
closure forces open-endedness, and the monotone count forces non-duplication
\textup{(}\Cref{thm:open,thm:nondup}\textup{)}.
\end{enumerate}
\end{proposition}

\begin{proof}
Each clause restates the cited theorem's use of \Cref{ax:finite}; the proofs are
those given in \Cref{sec:floor,sec:negation,sec:residual,sec:invariant,%
sec:gating,sec:monotone,sec:quotient,sec:construct}.
\end{proof}

\subsection{Failure without finiteness}

We show the ladder is not forced once the premise is dropped, by exhibiting, for
each load-bearing structure, a non-finite or zero-floor graph in which it fails.

\begin{proposition}[Necessity of the premise]
\label{prop:necessity}
Drop either clause of \Cref{ax:finite}:
\begin{enumerate}[label=\textup{(\alph*)},leftmargin=2.2em]
\item \textbf{Zero floor.} If edge weights may approach $0$
\textup{(}$\inf_e\wt(e)=0$\textup{)}, then a part may be separated from the rest
by a cut of arbitrarily small weight; $\wt(\cut(U))$ is no longer bounded below,
so T0 fails, and with it the residual positivity of T2 and the gate-tie
guarantees that rely on a floor.
\item \textbf{Infinite vertex set.} If $\abs{\V}=\infty$, a vertex may have
infinite degree, so a gate must break infinitely many ties at once and the finite
tie-breaking of T5 is not available in the same form; minimising partitions need
not be attained, so the residual of T2 need not be realised; and a walk's
``committed count'' may fail to be a well-ordered monotone natural number if
steps are not discrete, weakening T6.
\item \textbf{Complete graph with vanishing weights.} On an infinite complete
graph with weights tending to $0$, every part is adjacent to every other at
negligible cost: there is no positive residual, no privileged complement
structure, and no forced gate---none of T0--T8 is compelled.
\end{enumerate}
Hence the substructures are consequences of the premise, not of graphs in
general.
\end{proposition}

\begin{proof}
(a) Take a graph with a sequence of cuts of weight $2^{-n}\to0$ separating a
fixed part; then $\inf_U\wt(\cut(U))=0$, contradicting the conclusion of
\Cref{thm:floor}, which used $\wt(e)\ge\floorc>0$. (b) On an infinite star, the
centre has infinite degree, so no single sub-normalised distribution over its
incident edges selects after finitely many comparisons; and the partition lattice
of an infinite vertex set need not attain its infimum, so $\Res(\Gph)$ as a
minimum may not exist. (c) On $K_\infty$ with $\wt(e_n)=2^{-n}$, for any part $U$
one may route its separation through ever-cheaper edges so that
$\inf\wt(\cut(U))=0$; the floor, the positive residual, and the negation-based
individuation (which needs a determinate complement structure) all dissolve.
Each construction violates the conclusion of the corresponding rung, so finiteness
with a positive floor is necessary.
\end{proof}

\subsection{The Master Theorem}

\begin{theorem}[Structure is the shadow of finitude]
\label{thm:master}
Let $\Gph$ be a graph. The substructures T0--T8 are all forced if $\Gph$ is a
finite contact graph \textup{(}\Cref{ax:finite}\textup{)}, and none is forced if
$\Gph$ is infinite or admits edge weights with infimum $0$. Consequently a graph
is necessarily structured---carries the floor, the negation-determined parts, the
conserved region-valued residual, the private invariants, the gated trajectories,
the non-returning histories, the single-gated collectives, and the bounded
authorability---\emph{if and only if} it is finite with a positive floor. Structure
is neither imposed on a world from outside nor present in it independently of its
parts; it is exactly the form a world of finitely many positively separated parts
is compelled to take.
\end{theorem}

\begin{proof}
The forward direction is \Cref{prop:dependence} together with the theorems it
cites: under \Cref{ax:finite} each of T0--T8 holds. The reverse direction is
\Cref{prop:necessity}: dropping finiteness or the positive floor defeats each
load-bearing rung by explicit construction. Hence the conjunction of T0--T8 is
equivalent, over graphs, to the Finiteness Premise: finite-with-floor on one
side, the full ladder of structure on the other. The closing statement is the
biconditional read in words.
\end{proof}

\subsection{Scope}

The theorems are statements about finite weighted graphs and an embedded gating
datum; nothing more is assumed. The single premise (\Cref{ax:finite}) is a finite
graph whose edges cost something. We have made no claim about any particular
system. Where each rung uses the premise is marked (\Cref{prop:dependence}); that
the premise is necessary is shown by counterexample (\Cref{prop:necessity}).
Interpretive readings---of parts as agents, gates as selecting fields, quotient
gates as collective purpose, residuals as conserved truth, private invariants as
identities, authored structures as manufactured individuals---are confined to
remarks (\Cref{rem:negation-interp,rem:repacking,rem:character,rem:field,%
rem:purpose,rem:onelaw,rem:new-category}) on which no theorem depends. A reader
will recognise that finitely many bounded parts---of many kinds---satisfy the
premise; we prove everything of the abstract graph alone, and material that could
not be proved at this standard has been omitted.

\subsection{Conclusion}

From the single hypothesis that a world is a finite graph of positively separated
parts, an entire architecture follows of necessity: no separation is free; parts
are fixed by what they are not; the boundary content between parts is a conserved
region that no relabelling removes and no point realises; each part carries a
private conserved invariant beneath its changing trajectories; those trajectories
require a selector and never return to a prior state; a coherent collective of
parts has exactly one selector on its quotient; and such structures, being
completely specifiable, can be authored---though authorship reaches only the
instrument and not its music, produces only new individuals and never copies, and
resolves only to its own floor. Each is forced by finiteness and dissolves
without it. The answer to why a world is structured at all is therefore not that
structure is added to it or discovered in it, but that structure is what a world
\emph{must} be for finitely many bounded parts to stand apart within it at all.
Structure is the shadow of finitude; the contact graph is its form.

% =====================================================================
\bibliographystyle{plainnat}
\bibliography{references}
% =====================================================================

\end{document}
